%------------------------------------------------------------------------------%
\begin{exercises}

%------------------------------------------------------------------------------%
\begin{exercise}
  Nos itens que seguem, encontre o volume do sólido obtido pela rotação
  da região limitada pelas curvas dadas em torno das retas especificadas.
\begin{alphalist}
  \item $y=\sqrt[3]{x}$, $y=0$, $x=0$, $x=1$, em torno do eixo $x$.
  \item $y=x^2$,$y=6x-2x^2$, em torno do eixo $x$.
  \item $y=\frac{1}{x}$, $y=0$, $x=1$, $x=2$, em torno do eixo $x$.
  \item $y=x^2-6x+10$, $y=-x^2+6x-6$ em torno do eixo $x$.
  \item $y=\sqrt{x}$, $x=0$, $y=2$,  em torno do eixo $y$.
  \item $y=1+y^2$, $x=0$, $y=1$, $y=2$, em torno do eixo $y$.
  \item $y=x^3$, $y=8$, $x=0$ em torno do eixo $y$.
  \item $y=4x-x^2$, $y=3$, em torno do eixo $x=1$.
  \item $y=\sqrt{x}$, $y=0$, $x=1$ em torno do eixo $x=-1$.
  \item $y=x^3$, $y=0$, $x=1$ em torno do eixo $y=1$.
  \item $x=y^2+1$, $x=2$,  em torno do eixo $y=-1$.
\end{alphalist}
\begin{answer}
\begin{mathlist}[3]
  \item \frac{6\pi}{7}
  \item 8\pi
  \item \pi\left(1-\frac{1}{e}\right)
  \item \frac{16\pi}{3}
  \item 8\pi
  \item \frac{21\pi}{2}
  \item \frac{768\pi}{7}
  \item \frac{8\pi}{3}
  \item \frac{32\pi}{15}
  \item \frac{5\pi}{14}
  \item \frac{16\pi}{3}
\end{mathlist}
\end{answer}
\end{exercise}

%------------------------------------------------------------------------------%
\begin{exercise}
Nos itens a seguir, calcule o volume do sólido gerado pela rotação da região
limitada por $y=x$ e $y=2x-x^2$ em torno  das retas especificadas.
Você pode decidir qual dos dois métodos utilizados você deve usar.
\begin{alphalist}
  \item do eixo $x$
  \item do eixo $y$
  \item Da reta $x=1$
  \item da reta $y=2$
\end{alphalist}
%\begin{answer}
%\begin{mathlist}[2]
%  \item .
%  \item .
%  \item .
%  \item .
%\end{mathlist}
%\end{answer}
\end{exercise}

%------------------------------------------------------------------------------%
\begin{exercise}
Nos itens a seguir, calcule o volume do sólido gerado pela rotação do
triângulo com vértices $(1,1)$, $(1,2)$ e $(2,2)$em torno  das retas
especificadas.
Você pode decidir qual dos dois métodos utilizados você deve usar.
\begin{alphalist}
  \item do eixo $x$
  \item do eixo $y$
  \item Da reta $x=\frac{10}{3}$
  \item da reta $y=1$
\end{alphalist}
\begin{answer}
\begin{mathlist}[2]
  \item \frac{5\pi}{3}
  \item \frac{4\pi}{3}
  \item 2\pi
  \item \frac{2\pi}{3}
\end{mathlist}
\end{answer}
\end{exercise}

%------------------------------------------------------------------------------%
\begin{exercise}
  Seja $V$ o volume do sólido obtido pela rotação da região
  limitada por $y=\sqrt{x}$ e $y=x^2$ em torno do eixo $y$. Encontre $V$
  pelos métodos das seções transversais e das cascas cilíndricas.
  Em ambos os casos, desenhe um diagrama para explicar o método.
\begin{answer}
  $\frac{3\pi}{10}$
\end{answer}
\end{exercise}

%------------------------------------------------------------------------------%
\begin{exercise}
Use o método das cascas cilíndricas, encontre o volume do solido descrito.
\begin{alphalist}
  \item uma pirâmide de base quadrada com lado $L$ e cuja altura é $h$.
  \item Um cone circular reto de altura $h$ e raio da base $r$.
  \item Um tronco de cone circular reto de altura $h$, raio da
        base inferior $R$ e  raio da base superior $r$.
  \item Uma calota de uma esfera de raio $r$ e altura $h$.
\end{alphalist}
\begin{answer}
\begin{mathlist}[2]
  \item \frac{L^2h}{3}
  \item \frac{1}{3}r^2h
  \item \pi h \left( R^2 + Rh+ r^2 \right)
  \item \pi h^2 \left(r -\frac{h}{3} \right)
\end{mathlist}
\end{answer}
\end{exercise}

%------------------------------------------------------------------------------%
\begin{exercise}
  Escolha qual dos dois métodos estudados utilizar e encontre o volume
  do sólido gerado pela região dada em torno do eixo especificado.
\begin{alphalist}
  \item $y=-x^2+6x-8$, $y=0$ em torno do eixo $y$.
  \item $y=-x^2+6x-8$, $y=0$ em torno do eixo $x$.
\end{alphalist}
\begin{answer}
\begin{mathlist}[2]
 \item 8\pi
\end{mathlist}
\end{answer}
\end{exercise}

%------------------------------------------------------------------------------%
%\begin{exercise}
%\begin{mathlist}[2]
%  \item
%  \item
%  \item
%\end{mathlist}
%\begin{answer}
%  \begin{alphalist}[3]
%    \item
%    \item
%    \item
%  \end{alphalist}
%\end{answer}
%\end{exercise}

\end{exercises}
%------------------------------------------------------------------------------%
